% !TEX root = ./main.tex

\sysusetup{
  title                = {基于并行计算的金融大规模数据异常检测系统研究与实现},
  title*               = {Research and Implementation of a Parallel-Computing-Based Anomaly Detection System for Large-Scale Financial Data},
  author               = {黄裕文},
  author*              = {Yuwen Huang},
  speciality           = {工程硕士~人工智能（高性能计算方向）},
  speciality*          = {Professional Master of Engineering in Artificial Intelligence (HPC Track)},
  supervisor           = {黄聃~副教授},
  supervisor*          = {Dan Huang},
  % practice-supervisor  = {XXX~教授, XXX~教授},
  % practice-supervisor* = {Prof. XXX, Prof. XXX},
  % date                 = {2017-05-01},  % 完成时间，默认为今日
  % professional-type    = {工程硕士},
  % professional-type*   = {Master of Engineering},
   department           = {计算机学院},
  % student-id           = {[学号]},
  % secret-level         = {秘密},     % 绝密|机密|秘密|控阅，注释本行则公开
  % secret-level*        = {Secret},  % Top secret | Highly secret | Secret
  % secret-year          = {10},      % 保密/控阅期限
  % reviewer             = true,      % 声明页显示“评审专家签名”
  %
  % 数学字体
  % math-style           = GB,  % 可选：GB, TeX, ISO
  math-font            = xits,  % 可选：stix, xits, libertinus
}

% 模板已默认按学校要求使用“章-序号”格式编号，如“图2-1”“表2-1”“式（3-1）”。

% 加载宏包

% 定理类环境宏包
\usepackage{amsthm}

% 控制浮动体不跨节乱跑（保证图表按叙述顺序出现）
\usepackage[section]{placeins}

% 表注
\usepackage{threeparttable}

% 表格增强（列格式控制）
\usepackage{array}
\usepackage{ragged2e}

% 跨页表格
\usepackage{longtable}
\setlength{\LTleft}{0pt}
\setlength{\LTright}{0pt}

% 绘图（黑白结构图等）
\usepackage{tikz}
\usetikzlibrary{arrows.meta,positioning,calc}

% 算法
\usepackage[ruled,linesnumbered]{algorithm2e}
%\renewcommand{\thealgocf}{\thechapter-\arabic{algocf}}% 如需自定义算法编号，可在此基础上调整

% SI 量和单位
\usepackage{siunitx}

% 参考文献使用 BibTeX + natbib 宏包
% 顺序编码制
\usepackage[sort]{natbib}
\bibliographystyle{sysuthesis-numerical}

% 著者-出版年制
% \usepackage{natbib}
% \bibliographystyle{sysuthesis-authoryear}

% 参考文献使用 BibLaTeX 宏包
% \usepackage[style=sysuthesis-numeric]{biblatex}
% \usepackage[bibstyle=sysuthesis-numeric,citestyle=sysuthesis-inline]{biblatex}
% \usepackage[style=sysuthesis-authoryear]{biblatex}
% 声明 BibLaTeX 的数据库
% \addbibresource{reference.bib}

% 配置图片的默认目录
\graphicspath{{figures/}}

% 数学命令
\makeatletter
\newcommand\dif{%  % 微分符号
  \mathop{}\!%
  \ifsysu@math@style@TeX
    d%
  \else
    \mathrm{d}%
  \fi
}
\makeatother
\newcommand\eu{{\symup{e}}}
\newcommand\iu{{\symup{i}}}

% 用于写文档的命令
\DeclareRobustCommand\cs[1]{\texttt{\char`\\#1}}
\DeclareRobustCommand\env[1]{\texttt{#1}}
\DeclareRobustCommand\pkg[1]{\textsf{#1}}
\DeclareRobustCommand\file[1]{\nolinkurl{#1}}

\makeatletter
\renewcommand\sysu@notation@name{符号和缩略语说明}
\makeatother

\newcommand\thesisplaceholderbox[2][0.82\linewidth]{%
  \fbox{%
    \parbox[c][5cm][c]{#1}{%
      \centering #2%
    }%
  }%
}

% 统一表格列类型，中文英文长内容均可在单元格内自然换行
\newcolumntype{P}[1]{>{\RaggedRight\arraybackslash\hspace{0pt}}p{#1}}
\newcolumntype{Y}{>{\RaggedRight\arraybackslash\hspace{0pt}}X}
\newcolumntype{C}[1]{>{\Centering\arraybackslash\hspace{0pt}}p{#1}}

% 统一表格风格与版心约束
\captionsetup[table]{font=footnotesize, skip=6pt}
\AtBeginEnvironment{table}{%
  \footnotesize
  \setlength{\tabcolsep}{3.5pt}%
  \renewcommand{\arraystretch}{1.18}%
}
\AtBeginEnvironment{tabular}{%
  \footnotesize
  \setlength{\tabcolsep}{3.5pt}%
  \renewcommand{\arraystretch}{1.18}%
  \setlength{\emergencystretch}{1em}%
}
\AtBeginEnvironment{tabularx}{%
  \footnotesize
  \setlength{\tabcolsep}{3.5pt}%
  \renewcommand{\arraystretch}{1.18}%
  \setlength{\emergencystretch}{1em}%
}
\AtBeginEnvironment{longtable}{%
  \footnotesize
  \setlength{\tabcolsep}{3.5pt}%
  \renewcommand{\arraystretch}{1.18}%
  \setlength{\emergencystretch}{1em}%
}

% hyperref 宏包在最后调用
\usepackage{hyperref}
